\section{Verträge im Netz}

\subsection{Vertragsschluss und Konsensus im E-Commerce}

\begin{definition}
    \textbf{Ablauf des Vertragsschlusses im Webshop:} Ein Vertrag entsteht durch gegenseitige, übereinstimmende Willensäusserung (Konsens, Art. 1 OR):
    \begin{enumerate}
        \item \textbf{Einladung zur Offerte (\textit{invitatio ad offerendum}):} Die Präsentation von Waren oder Dienstleistungen auf einer Website oder in einem Webshop stellt in der Regel noch keinen rechtlich verbindlichen Antrag dar (Art. 7 Abs. 2 OR), sondern lediglich eine unverbindliche Aufforderung an den Kunden, seinerseits eine Offerte abzugeben.
        \item \textbf{Angebot (Bestellung):} Durch das Anklicken des Bestell-Buttons (z.B. ,,Zahlungspflichtig bestellen``) gibt der Kunde ein verbindliches Angebot zum Abschluss eines Kaufvertrags ab.
        \item \textbf{Annahme:} Der Vertrag kommt zustande, sobald der Händler das Angebot annimmt. Dies geschieht entweder durch eine ausdrückliche Annahmeerklärung (z.B. Auftragsbestätigung, Versandbestätigung per E-Mail) oder konkludent durch die tatsächliche Lieferung der Ware.
    \end{enumerate}
\end{definition}

\begin{important}
    \textbf{Automatische Empfangsbestätigung $\neq$ Vertragsannahme:} Die gesetzlich vorgeschriebene sofortige elektronische Bestätigung des Bestelleingangs (Art. 3 Abs. 1 lit. s Ziff. 4 UWG) dient nur der Dokumentation des Eingangs und stellt noch keine Annahme der Offerte dar, es sei denn, der Text erklärt dies ausdrücklich.
\end{important}

\begin{theorem}
    \textbf{Willensmängel und Irrtümer im E-Commerce (Art. 23 ff. OR):}
    Kommt es bei der Erfassung oder automatisierten Verarbeitung im Webshop zu Preis- oder Produktfehlern, stellt sich die Frage der Vertragsgültigkeit:
    \begin{itemize}
        \item \textbf{Wesentlicher Erklärungsirrtum (Art. 24 Abs. 1 Ziff. 3 OR):} Weicht die abgegebene Erklärung unbewusst vom tatsächlichen Willen ab (z.B. Tippfehler bei der Dateneingabe, Datenbank-Synchronisationsfehler), liegt ein wesentlicher Erklärungsirrtum vor, wenn eine Leistung von erheblich grösserem oder geringerem Wert zugesagt wurde als beabsichtigt.
        \item \textbf{Rechtsfolge (Anfechtbarkeit):} Der Vertrag ist für den Irrenden unverbindlich, sofern er den Irrtum unverzüglich nach Entdeckung geltend macht (Art. 31 OR).
        \item \textbf{Fahrlässiger Irrtum und Schadenersatz (Art. 26 OR):} Beruht der Erklärungsirrtum auf eigener Fahrlässigkeit des Händlers, ist die Anfechtung dennoch wirksam; der Händler muss dem gutgläubigen Käufer jedoch den aus dem Dahinfallen des Vertrages erwachsenen Schaden (negatives Vertragsinteresse) ersetzen.
        \item \textbf{Praxisfall (Falsche Preisanschrift):} Bei offensichtlichen Preisfehlern (z.B. Microspot- oder Orell Füssli-Fall: Luxusartikel oder Bücher für CHF 0.-- oder ein Bruchteil des Marktpreises) fehlt es am Schutz des Käufers, wenn dieser nach Treu und Glauben erkennen musste, dass es sich um ein Versehen handelte (Rechtsmissbrauchsverbot, Art. 2 Abs. 2 ZGB).
    \end{itemize}
\end{theorem}

\begin{definition}
    \textbf{Vertragsschluss durch KI-Agenten und Bots:}
    Immer häufiger werden Verträge ganz oder teilweise durch automatisierte Systeme oder generative KI-Agenten initiiert, ausgehandelt oder abgeschlossen:
    \begin{itemize}
        \item \textbf{Rechtliche Zurechnung:} KI-Systeme besitzen keine eigene Rechtspersönlichkeit. Erklärungen, die ein KI-Agent abgibt, werden dem Betreiber (Unternehmen) nach den Grundsätzen des Erklärungsboten oder der Stellvertretung analog zugerechnet: Der Betreiber bedient sich des Systems als Werkzeug zur Willensäusserung.
        \item \textbf{Verbindlichkeit und Risiken:} Äussert ein Chatbot gegenüber einem Kunden Zusagen, Rabatte oder Vertragsbedingungen (z.B. durch Fehlinterpretationen oder Halluzinationen wie im Fall des Air-Canada-Chatbots), ist das Unternehmen grundsätzlich an diese Zusagen gebunden.
        \item \textbf{Schutzmöglichkeiten des Betreibers:}
        \begin{itemize}
            \item Technische Restriktionen (Hard Limits bei Rabatten, strikte System-Prompts, Freigabevorbehalte durch Menschen bei kritischen Abschlüssen).
            \item Transparenzhinweis, dass ein KI-System agiert.
            \item Einbeziehung klarer AGB-Klauseln, wonach KI-Auskünfte unverbindlich sind und ein Vertrag erst mit menschlicher Bestätigung zustande kommt.
        \end{itemize}
    \end{itemize}
\end{definition}

\subsection{Allgemeine Geschäftsbedingungen (AGB)}

\begin{definition}
    \textbf{Einbezug von AGB in den Vertrag:} AGB sind vorformulierte Vertragsbestimmungen für eine Vielzahl von Verträgen. Sie werden nur dann Vertragsbestandteil, wenn sie wirksam einbezogen wurden:
    \begin{itemize}
        \item \textbf{Vorheriger deutlicher Hinweis:} Der Kunde muss \textbf{vor} Abgabe seiner Vertragserklärung ausdrücklich auf die Geltung der AGB hingewiesen werden.
        \item \textbf{Zumutbare Kenntnisnahmemöglichkeit:} Dem Kunden muss es vor Vertragsschluss zumutbar möglich sein, vom Inhalt der AGB Kenntnis zu nehmen (z.B. über einen gut sichtbaren, funktionsfähigen Hyperlink direkt neben der Bestell-Checkbox).
        \item \textbf{Zustimmung (Opt-in):} Bestätigung durch aktives Anklicken einer Checkbox (,,Ich akzeptiere die AGB``). Vorauszugefüllte Checkboxen oder blosse Hinweise im Footer genügen für eine wirksame Einbindung oft nicht.
    \end{itemize}
\end{definition}

\begin{theorem}
    \textbf{Inhaltskontrolle von AGB:}
    Nicht jede wirksam einbezogene Klausel ist auch rechtlich durchsetzbar:
    \begin{itemize}
        \item \textbf{Ungewöhnlichkeitsregel (Praxis):} Eine Klausel gilt als nicht vereinbart, wenn sie geschäftsfremden Charakter aufweist, der Kunde ihr nicht ausdrücklich zugestimmt hat und nach Treu und Glauben vernünftigerweise nicht mit einer solchen Bestimmung rechnen musste.
        \item \textbf{Gesetzliche Inhaltskontrolle (Art. 8 UWG):} Im Verhältnis zu Konsumentinnen und Konsumenten (B2C) sind AGB-Klauseln unlauter und damit \textbf{nichtig}, die in Treu und Glauben verletzender Weise ein erhebliches und ungerechtfertigtes Missverhältnis zwischen den vertraglichen Rechten und Pflichten zum Nachteil des Konsumenten begründen. \\
        Typische unzulässige Klauseln: Umfassende Haftungsausschlüsse für Kernpflichten, einseitige nachträgliche Preisänderungsrechte ohne Kündigungsoption, unzumutbare Gerichtsstandsklauseln.
    \end{itemize}
\end{theorem}

\subsection{Informationspflichten und Konsumentenschutz}

\begin{definition}
    \textbf{Gesetzliche Informationspflichten im E-Commerce (Art. 3 Abs. 1 lit. s UWG):}
    Wer Waren, Werke oder Dienstleistungen im elektronischen Geschäftsverkehr anbietet, muss:
    \begin{enumerate}
        \item \textbf{Impressumspflicht:} Klare und vollständige Angaben über seine Identität und seine Kontaktadresse (einschliesslich E-Mail) machen.
        \item \textbf{Transparenz der Schritte:} Auf die einzelnen technischen Schritte hinweisen, die zum Vertragsschluss führen.
        \item \textbf{Korrekturmöglichkeiten:} Angemessene technische Mittel zur Verfügung stellen, mit denen Eingabefehler vor dem Absenden der Bestellung erkannt und korrigiert werden können (z.B. Bestellübersichtsseite).
        \item \textbf{Unverzügliche Bestätigung:} Die Bestellung des Kunden unverzüglich auf elektronischem Weg (z.B. Bestätigungs-E-Mail) bestätigen.
    \end{enumerate}
\end{definition}

\begin{important}
    \textbf{Kein gesetzliches Widerrufsrecht im Schweizer Online-Handel (Unterschied CH vs. EU):}
    \begin{itemize}
        \item \textbf{EU-Recht:} Im europäischen Fernabsatzrecht (EU-Verbraucherrechte-Richtlinie) gilt ein zwingendes, mindestens \textbf{14-tägiges gesetzliches Widerrufsrecht} für Verbraucher bei Online-Käufen ohne Angabe von Gründen.
        \item \textbf{Schweizer Recht:} Das schweizerische Recht kennt \textbf{kein allgemeines Widerrufs- oder Rückgaberecht} für Online-Käufe! Das Widerrufsrecht für Konsumverträge gemäss Art. 40a ff. OR (14 Tage) gilt ausschliesslich für Haustürgeschäfte, Verkaufsfahrten und Telefonmarketing; Verträge im Webshop sind davon ausdrücklich ausgenommen. Ein Rückgaberecht im Schweizer Online-Handel existiert nur, wenn der Händler es freiwillig (z.B. in den AGB) gewährt.
    \end{itemize}
\end{important}

\subsection{Beweis, Schriftform und Elektronische Signatur}

\begin{definition}
    \textbf{Formvorschriften im Schweizer Recht (Art. 11 ff. OR):}
    \begin{itemize}
        \item \textbf{Grundsatz der Formfreiheit (Art. 11 Abs. 1 OR):} Verträge bedürfen zu ihrer Gültigkeit nur dann einer besonderen Form, wenn das Gesetz dies vorschreibt oder die Parteien dies vereinbart haben. Die meisten E-Commerce-Kaufverträge sind formfrei und per Mausklick gültig.
        \item \textbf{Einfache Schriftlichkeit (Art. 13/14 OR):} Wo das Gesetz Schriftform verlangt (z.B. Lehrvertrag, Schenkungsversprechen, Zession Art. 165 OR, Konsumkreditvertrag), muss das Dokument eigenhändig von der verpflichteten Person unterschrieben werden. Eine blosse E-Mail oder ein Scan genügen diesem Erfordernis nicht.
    \end{itemize}
\end{definition}

\begin{theorem}
    \textbf{Qualifizierte elektronische Signatur (Art. 14 Abs. 2bis OR i.V.m. ZertES):}
    Der eigenhändigen Unterschrift ist rechtlich ausschliesslich die \textbf{qualifizierte elektronische Signatur (QES)} gleichgestellt, die auf einem qualifizierten Zertifikat einer anerkannten Schweizer Zertifizierungsstelle beruht und mit einem qualifizierten elektronischen Zeitstempel versehen ist (gemäss Bundesgesetz über die elektronische Signatur, ZertES). Nur die QES erfüllt gesetzliche Schriftformerfordernisse digital.
\end{theorem}

\begin{remark}
    \textbf{Sorgfalts- und Haftungspflichten des Signaturinhabers (Art. 59a OR):}
    Der Inhaber eines Signaturschlüssels ist gesetzlich verpflichtet, die Signierschlüsseldaten sorgfältig zu verwahren und vor Missbrauch durch Unbefugte zu schützen. Verletzt er diese Sorgfaltspflicht (z.B. Weitergabe der Signatur-PIN), haftet er Dritten für den daraus entstandenen Schaden, wenn diese auf die Gültigkeit der Signatur vertraut haben.
\end{remark}